\chapter{Вовед, мотивација и теоретска основа}
\label{chap:intro}

\section{Проблем и мотивација}

Четириножниот робот може да биде стабилен и функционален, а сепак неговото
однесување да биде тешко читливо за човек. При разговорна интеракција не е
доволно системот само да генерира текстуален одговор: корисникот очекува
телото, ритамот и паузата на роботот да бидат усогласени со содржината. Во
афективното пресметување емоцијата не се третира само како декоративен слој,
туку како информација што учествува во перцепција, одлучување и интеракција.
Кај социјалните роботи, експресивното однесување е
дел од координацијата со човекот и мора да се оценува и како изолиран сигнал и
во целосната интеракција.

Проблемот е потежок кај Lite3 од две причини. Прво, роботот нема човечко лице,
раце или артикулирана глава; значењето мора да се пренесе преку висина на
трупот, roll, pitch, темпо, симетрија, пауза и контролирано движење на
стапалата. Второ, истата идеја не смее директно да се копира од Gazebo на
физичкиот робот. Симулацискиот контролер прима апстрактна Joy команда и има
посебни hop/stomp механизми, додека физичкиот MotionSDK испраќа 12 joint
траектории под услови на контакт, свеж feedback и ексклузивна сопственост.

Оттука е поставено главното проектно прашање: \emph{како разговорно
добиената емоционална состојба да се претвори во читливо движење, без
емоционалното расудување да добие директен и неограничен пристап до актуаторите?}

\section{Цели}

Во трудот се поставени следните цели:

\begin{enumerate}[label=\arabic*.]
  \item да се интегрира headless EmotionBot со ROS Noetic без GUI, микрофон,
  network или model-loading side effects при import;
  \item да се дефинираат верзионирани, проверливи договори за разговорен turn,
  емоционална состојба и expression action;
  \item да се изведат девет визуелно различни, центрирани Gazebo профили;
  \item да се воведе повеќеслојна симулациска заштита со clamp, arbitration,
  readiness, watchdog и точна нула при stale/disable/shutdown;
  \item да се изгради одделна split-host физичка архитектура со еден MotionSDK
  owner и torque-derived contact gates;
  \item да се документираат позитивните и негативните експерименти без да се
  мешаат доказните нивоа; и
  \item да се обезбеди репродуцибилен Docker/Git workflow.
\end{enumerate}

\section{Афективно и социјално роботско однесување}

Социјално интерактивниот робот не се дефинира само според механичката форма,
туку и според улогата на социјалната интеракција, употребените знаци и
очекувањата што тие ги создаваат.
Движењето во овој проект затоа се третира како комуникациски канал. Забрзаното
подигање и freeze-паузата можат да упатуваат на изненадување; бавното
спуштање на предниот дел на телото може да упатува на тага; наизменичното
подигање на предните шепи може да се прочита како возбуда. Овие интерпретации
не се универзална психолошка вистина, туку дизајнерски хипотези што мора да се
проверат визуелно од оператор и да останат во механички безбедна обвивка.

Во оваа смисла, „емоција“ во системот има три нивоа:

\begin{enumerate}
  \item \textbf{внатрешна состојба} --- valence, arousal и категорија што ги
  одржува EmotionEngine;
  \item \textbf{комуникациска намера} --- декларативен профил или action што
  ROS-слојот го извлекува од состојбата; и
  \item \textbf{изведено однесување} --- симулациско или физичко движење што
  ги поминало релевантните gates.
\end{enumerate}

Разделувањето спречува јазичниот модел или категоризаторот директно да задава
joint, torque или UDP вредности.

\section{Valence--arousal модел и дискретни категории}

Циркумплексниот модел го опишува афектот со две приближно ортогонални димензии:
пријатност/непријатност и активација. Во проектот
се користат нормализирани величини

\begin{equation}
  v \in [-1,1], \qquad a \in [0,1],
  \label{eq:va-bounds}
\end{equation}

каде што $v$ е valence, а $a$ е arousal. Дискретната категорија
$e$ е член на множеството

\begin{equation}
  \mathcal{E}=\{\text{neutral, joy, sadness, anger, fear, surprise, disgust,
  curiosity, affection}\}.
  \label{eq:emotion-set}
\end{equation}

Двете претставувања се комплементарни. Категоријата ја избира кореографијата,
а valence/arousal ја менуваат интензивноста во однапред ограничен интервал.
Системот не дозволува висок arousal да ги прошири hard safety bounds. Со тоа
психолошкиот модел служи за изразување, а не за авторизација.

\section{Движење како емоционален израз}

Кај робот без лице се користат четири главни параметри: просторна насока,
амплитуда, динамика и повторување. Брза промена со висока положба се чита
поинаку од бавен, низок и асиметричен став. Сепак, препознатливоста не може да
се заклучи само од вредности во лог. Во овој проект физички профил е
„прифатен“ дури кога: (1) телеметријата ги потврдила gates, (2) движењето
завршило и сопственоста била ослободена и (3) операторот визуелно го оценил
како соодветна емоција.

Ова правило објаснува зошто безбедни, но премногу суптилни Joy и Sadness
прототипи се заведени како неуспешни. И обратно, позитивна визуелна оценка не
го поништува telemetry failure: првиот 15-секунден Anger repeat изгледал добро,
но третиот циклус завршил со само три повторно потврдени supports и затоа не е
евидентиран како успешна endurance проверка.

\section{ROS и симулацијата како метод}

ROS обезбедува процесно разделување преку nodes, topics, services и parameter
server. Gazebo овозможува динамички модел,
контакти и контролери за повторливи експерименти пред физичка активација. Во
овој проект симулацијата има три
улоги: рано откривање contract/ordering грешки, визуелно обликување на деветте
профили и автоматизирана проверка на bounds, zeroing и no-hardware boundary.

Симулацијата не е доказ за физичка безбедност. Моделот, контактниот solver,
актуаторите и transport path се различни. Затоа архитектурата намерно се
разгранува по емоционалната состојба, а не по веќе пресметана симулациска
команда. Оваа одлука е основа на целиот дизајн.
